\begin{helper}[Product Tail Comparison]
Fix \(n\in\mathbb N\), \(I\subseteq\mathrm{Fin}(n)\), parameters
\(\varepsilon,\varepsilon_1,\varepsilon_2\in\mathbb R\), and
\(n_0\in\mathbb N\).  Assume
\[
\noderef{ParameterHierarchyEventualComparisons}
  (\varepsilon,\varepsilon_1,\varepsilon_2,n_0),\qquad n_0<n.
\]
Let
\[
M=\noderef{TwoBiteNaturalReducedVertexCount}(n),\qquad
p_0=\noderef{TwoBiteEdgeProbability}\!\left(\frac12,n\right).
\]
For an embedding
\[
\pi:\mathrm{Fin}(n)\hookrightarrow\mathrm{Fin}(M)\times\mathrm{Fin}(M),
\]
put
\[
P_R=\{\pi(u)_1:u\in I\},\qquad P_B=\{\pi(u)_2:u\in I\},
\]
\[
q=\noderef{FixedSetProductModelMassFn}(p_0,P_R,P_B),
\qquad
Z_{\mathrm{prod}}=\noderef{FixedSetProductModelVar}(I,\varepsilon,p_0,\pi).
\]
Let \(t,\mu\in\mathbb R\).  Assume the product mean satisfies
\[
\mathbb E_q Z_{\mathrm{prod}}
=\sum_v\left(\prod_i q(v_i)\right)Z_{\mathrm{prod}}(v)
\le \mu.
\]
Define \(Z_I\) from a two-bite sample by the same external-small-class,
final-pair, and distinct-projection count used in
\(\noderef{FixedSetProductModelVar}\).  Then
\[
\begin{aligned}
&\noderef{TwoBiteEventProbability}(n,M,p_0)
 \bigl(\{\omega:\noderef{TwoBiteEmbedding}(\omega)=\pi
   \text{ and } Z_I(\omega)\ge \mu+t\}\bigr)\\
&\quad\le
\Pr_q\!\left(Z_{\mathrm{prod}}\ge
  \mathbb E_qZ_{\mathrm{prod}}+t\right)
\noderef{TwoBiteEventProbability}(n,M,p_0)
 \bigl(\{\omega:\noderef{TwoBiteEmbedding}(\omega)=\pi\}\bigr),
\end{aligned}
\]
where the product-tail probability is the finite sum of
\(\prod_iq(v_i)\) over assignments \(v\) satisfying
\[
\mathbb E_qZ_{\mathrm{prod}}+t\le Z_{\mathrm{prod}}(v).
\]
\end{helper}
\begin{proof}
Let
\[
\mathcal C_\pi
=\{\omega:\noderef{TwoBiteEmbedding}(\omega)=\pi\},
\qquad
E=\sum_v\left(\prod_iq(v_i)\right)Z_{\mathrm{prod}}(v).
\]
The hierarchy comparison and the identity
\(p_0=\noderef{TwoBiteEdgeProbability}(1/2,n)\) give
\[
0\le p_0\le1.
\]
Consequently every edge factor in \(\noderef{GnpGraphWeight}(p_0,G)\) is
nonnegative: an edge contributes \(p_0\), and a non-edge contributes
\(1-p_0\).  The injection factor
\(\noderef{UniformInjectionWeight}(n,M)\) is either \(0\) or the reciprocal of
a positive finite cardinal, hence is nonnegative.  Thus every
\(\noderef{TwoBiteSampleWeight}\) is nonnegative.  Also
\(\noderef{FixedSetProductModelMass}\), applied to \(P_R,P_B\), gives
\[
q(a)\ge0\quad\text{for all }a,\qquad \sum_a q(a)=1,
\]
so all product weights \(\prod_iq(v_i)\) and all product-tail sums are
nonnegative.  These nonnegativity facts justify every monotonicity step below.

First consider the case
\[
\noderef{TwoBiteEventProbability}(n,M,p_0)(\mathcal C_\pi)=0.
\]
The event
\(\mathcal C_\pi\cap\{Z_I\ge\mu+t\}\) is a subset of \(\mathcal C_\pi\).
Because all sample weights are nonnegative, its mass is at most the mass of
\(\mathcal C_\pi\), hence is \(0\).  The right side of the claimed inequality
is the nonnegative product-tail sum multiplied by this zero cell mass, so it
is also \(0\).  The desired inequality follows.

Assume from now on that the cell mass is positive.  On \(\mathcal C_\pi\), the
embedding is fixed, so the red and blue projections of vertices in \(I\) are
fixed as \(\pi(u)_1\) and \(\pi(u)_2\).  Hence the sets
\[
P_R=\{\pi(u)_1:u\in I\},\qquad P_B=\{\pi(u)_2:u\in I\},
\]
and the excluded embedded base set
\((P_R\text{ in the red copy})\cup(P_B\text{ in the blue copy})\)
are fixed on the cell.  Unfolding
\(\noderef{TwoBiteEventProbability}\) and \(\noderef{TwoBiteSampleWeight}\),
the weight of a sample \(\omega=(G_R,G_B,\pi)\in\mathcal C_\pi\) is
\[
\noderef{GnpGraphWeight}(p_0,G_R)\,
\noderef{GnpGraphWeight}(p_0,G_B)\,
\noderef{UniformInjectionWeight}(n,M).
\]
The injection factor is constant throughout the cell, so all conditional
comparisons inside the cell reduce to finite sums over the red and blue graph
edge variables.

Partition the red graph edge variables into three finite classes:
edges with both endpoints in \(P_R\), edges with both endpoints outside
\(P_R\), and incident edges \(\{x,r\}\) with \(x\notin P_R\) and
\(r\in P_R\).  Only the incident class can affect
\(\noderef{TwoBiteX}(\omega,I,\operatorname{inl} x)\), because every
vertex of \(I\) has red projection in \(P_R\).  The internal and disjoint
red edges are ignored by \(Z_I\); summing over each ignored edge contributes
\[
(1-p_0)+p_0=1.
\]
Thus, after summing out ignored red edges, a fixed family
\[
A_R(x)\subseteq P_R\qquad(x\in\mathrm{Fin}(M)\setminus P_R)
\]
has red weight
\[
\prod_{x\notin P_R}
p_0^{|A_R(x)|}(1-p_0)^{|P_R|-|A_R(x)|},
\]
where \(A_R(x)\) records exactly those \(r\in P_R\) adjacent to \(x\) in
\(G_R\).  The same partition for the blue graph gives incident data
\[
A_B(y)\subseteq P_B\qquad(y\in\mathrm{Fin}(M)\setminus P_B)
\]
with blue weight
\[
\prod_{y\notin P_B}
p_0^{|A_B(y)|}(1-p_0)^{|P_B|-|A_B(y)|},
\]
and all blue edges internal to \(P_B\) or disjoint from \(P_B\) sum out to
\(1\).

Now compare this conditional graph law with the product assignment law in
\(\noderef{FixedSetProductModelVar}\).  A product coordinate is a pair
\((R,B)\) of subsets.  For a red coordinate indexed by
\(x\notin P_R\), the fake red graph uses only the first component
\(R\), and the event data set \(A_R(x)\) is \(R\).  The second component is
redundant, and by \(\noderef{FixedSetProductModelMass}\)
\[
\sum_{B\subseteq P_B}
p_0^{|R|}(1-p_0)^{|P_R|-|R|}
\cdot
p_0^{|B|}(1-p_0)^{|P_B|-|B|}
=p_0^{|R|}(1-p_0)^{|P_R|-|R|};
\]
terms with \(R\not\subseteq P_R\) or \(B\not\subseteq P_B\) are zero by
\(\noderef{FixedSetProductModelMassFn}\).  If a red coordinate is indexed by
\(x\in P_R\), then \(\noderef{FixedSetFakeGraph}\) never reads the assignment
at \(x\), because \(x\) is itself on the embedded side and is removed from the
external base set; summing the whole pair \((R,B)\) over \(q\) gives \(1\) by
\(\noderef{FixedSetProductModelMass}\).

For a blue coordinate indexed by \(y\notin P_B\), the fake blue graph uses
only the second component \(B=A_B(y)\); the first component is redundant and
sums to \(1\) by the same calculation.  If \(y\in P_B\), the entire product
coordinate is ignored and again has total \(q\)-mass \(1\).  Therefore, after
all redundant product coordinates and redundant pair components are summed
out, the remaining product assignment data are in bijection with the incident
edge data
\[
\bigl(A_R(x)\bigr)_{x\notin P_R},
\qquad
\bigl(A_B(y)\bigr)_{y\notin P_B}.
\]
Under this bijection, the product weight \(\prod_iq(v_i)\) pushes forward to
exactly the conditional two-bite graph weight on the cell \(\mathcal C_\pi\).

For a sample in the cell and a product assignment with the same incident data,
the red graph used by \(\noderef{FixedSetProductModelVar}\) has precisely the
same adjacencies from outside \(P_R\) into \(P_R\) as \(G_R\), and the blue
fake graph has precisely the same adjacencies from outside \(P_B\) into
\(P_B\) as \(G_B\).  Since \(Z_I\) only uses
\(\noderef{TwoBiteSmallClass}\), \(\noderef{TwoBiteX}\),
\(\noderef{TwoBiteFinalPairs}\), and the red and blue projections of vertices
of \(I\), these incident adjacencies determine the value of \(Z_I\) on the
cell.  Thus the corresponding product assignment has
\[
Z_{\mathrm{prod}}(v)=Z_I(\omega).
\]
Equivalently, for every predicate \(H\) depending only on this value, finite
grouping by the incident data gives the pushforward identity
\[
\noderef{TwoBiteEventProbability}(n,M,p_0)
  (\mathcal C_\pi\cap\{\omega:H(Z_I(\omega))\})
=
\noderef{TwoBiteEventProbability}(n,M,p_0)(\mathcal C_\pi)
  \sum_{\{v:H(Z_{\mathrm{prod}}(v))\}}\prod_iq(v_i).
\]

Apply this identity to the tail comparison.  By the assumed mean bound,
\[
E\le\mu.
\]
Therefore every assignment satisfying
\[
\mu+t\le Z_{\mathrm{prod}}(v)
\]
also satisfies
\[
E+t\le Z_{\mathrm{prod}}(v).
\]
Using \(Z_I(\omega)=Z_{\mathrm{prod}}(v)\) under the pushforward
correspondence, the cell event \(Z_I\ge\mu+t\) is contained in the product
tail event \(E+t\le Z_{\mathrm{prod}}\).  The product-tail probability is
exactly the finite sum displayed in the theorem statement:
\[
\sum_{\{v:E+t\le Z_{\mathrm{prod}}(v)\}}\prod_iq(v_i).
\]
Multiplying this product-tail bound by the positive cell mass
\(\noderef{TwoBiteEventProbability}(n,M,p_0)(\mathcal C_\pi)\) gives the
claimed inequality for
\(\mathcal C_\pi\cap\{Z_I\ge\mu+t\}\).
\end{proof}
